\documentclass{article}
\usepackage{amsmath,hyperref}
\begin{document}
\section*{An aperture promise with independently scheduled leaf controls}
\paragraph{Source and modeled interface.}
The replay executes the selected pinned Open-RMF joint target,
command, mode, publication and fleet-phase functions with inert
Gazebo/ROS types. At the initial read the sampled positions are
fully open and the active session matches. The histories vary the
retained closing command before that read. Independent one-joint
models exercise the actual per-joint recurrence; giving distinct
leaf writers independent targets for one shared opening is a
declared design here, not a stock per-leaf interface to one door.

The geometry is an ideal opposed two-slider aperture. Each leaf's
open displacement is normalized to 1; usable aperture is the sum
of the two displacements, and the stipulated active traversal must
have a minimum gap at all times over 24 updates. Define the
nonnegative deficit relative to full aperture in units of $1/1024$.
The model applies the issued velocity immediately for each fixed
$1/8$-second interval, with an open stop. The closed stop never
binds in the sweep. Within an update, each deficit is a maximum of
zero and an affine function; their sum is convex and its interval
maximum is at an endpoint. There is no claim about inertial
hardware, collision avoidance outside this aperture condition, or
a Gazebo run. A separate entity's closing command executes in every
profile fixture and does not enter the aperture measure.

\paragraph{Exact pointwise sufficient state.}
Retain the no-pulse and proposed-pulse profiles for each leaf with
their clock alignment. A scope admits a switch of the target for
the indicated pulse, without requiring the optional maintenance
request actually to run. The envelope includes every combination
of the permitted occurrences. If $D_i(t,e_i)$ is the deficit of
leaf $i$ at cut $t$, then safety for every permitted occurrence
is equivalent to
\[
  \forall t,\quad \sum_{i=1}^{2}\max_{e_i\text{ permitted}}D_i(t,e_i)\le S(t).
\]
This follows by maximizing at each common cut over the product of
the two occurrence sets; the maximizing choices need not agree
between cuts. If the occurrence sets are coupled, the product can
be a safe overapproximation but not an exact criterion. Checking
the same capacity index at each cut reuses the admitted-row
checker. Direct enumeration of both actual leaf profiles ties.
Summing separate peaks is a sound, strictly coarser reference.
Holding both actuators is no remedy for an already excessive
retained stopping displacement, and excludes every optional pulse
even when the exact shared rule permits it.

\paragraph{Replacement and necessity.}
At 44 deficit units, the staggered pulses can both be allowed even
though their separately retained peaks sum to 76. At 162 units
in another matched fixture, either pulse by itself is compatible
with the other leaf's residual action, yet both cannot be allowed:
the aligned maximum is 316. There is no valid compatibility claim
from two separately admissible policies.
For one replacement, keep the left pulse at tick 9 and move the
right pulse from tick 5 to tick 11. The two standalone peaks remain
32 and 20. The joint maximum increases from 36 to 52, invalidating
the accepted 36-unit aperture promise. Pointwise domination at
the same clock would have preserved it. Lean checks the product
law, its pointwise replacement implication and the separate-peak
sufficiency, not the source-to-plant interpretation.

As in the digital revocation case, a fresh correct read does not
confer validity through an effect. Here the protected support is a
time-indexed geometric capacity, not a serialized permission row:
disjoint writers may still interfere, and a timed benign change
need not be fenced. Profile availability and the ability to enforce
admitted target schedules through the traversal are premises of
this held-in interpretation, not capacities supplied by the OPEN
report or demonstrated by an independent mission.
\end{document}
